\section{Alcance de implementación}

El proyecto contempla la implementación y parametrización de la plataforma Odoo Online (SaaS), seleccionada por su capacidad para centralizar operaciones en la nube sin requerir infraestructura local. El alcance se centra en funcionalidades estándar para asegurar una adopción ágil y segura de los siguientes componentes.

\subsection{Alcance temático}

Se han seleccionado estratégicamente los módulos que cubrirán la operatividad crítica de Machu Picchu Exclusive Tours:

\begin{itemize}
    \item[a.] \textbf{CRM turístico (gestión de clientes)}: centralización de la base de datos de pasajeros (PAX). Este módulo es fundamental para el almacenamiento seguro de documentos de identidad (pasaportes), eliminando el riesgo de seguridad que supone gestionarlos por WhatsApp. Incluye el seguimiento de oportunidades mediante un embudo de ventas (pipeline) y la segmentación básica con etiquetas.
    \item[b.] \textbf{Módulo de Proyectos (control operativo)}: configuración de un tablero visual (Kanban) para el seguimiento de la prestación del servicio. El flujo se adaptará a etapas como: Borrador de Cotización, Confirmación (Adelanto), Compra de Tickets/Hoteles y Ejecución final, permitiendo controlar itinerarios que varían de 2 a 20 días.
    \item[c.] \textbf{Gestión de cobros y saldos (facturación básica)}: uso del módulo estándar para el registro manual de adelantos (30\% a 50\%) y saldos que el cliente paga a través de enlaces externos (Ip-App). Esto permitirá visualizar el estado de pago (Pagado/Pendiente) para habilitar la compra de boletos no reembolsables con seguridad.
    \item[d.] \textbf{Gestión de actividades}: uso de recordatorios y alertas manuales (llamadas, correos) para el seguimiento post-cotización y la atención personalizada que caracteriza a la empresa.
\end{itemize}

Este alcance temático se encuentra alineado con el documento de Alcance del Proyecto, en el que se explicita que la implementación se LIMITA a las funcionalidades estándar de Odoo para CRM turístico, control de proyectos, actividades básicas y facturación simple, dejando fuera integraciones automáticas con portales de proveedores, módulos contables avanzados, e-commerce o desarrollo de reportes de inteligencia de negocios fuera de las vistas estándar. De este modo, se garantiza una implementación acotada, realista y coherente con el presupuesto y el horizonte temporal definidos.
\\
Al delimitar con precisión qué procesos quedan cubiertos y cuáles deberán abordarse en etapas posteriores, el alcance temático permite que la empresa enfoque sus esfuerzos en resolver primero los problemas más críticos: la dispersión de la información de clientes, la falta de trazabilidad de itinerarios y la ausencia de un control básico sobre adelantos y saldos. A partir de esta base, será posible evaluar, en futuras fases, la incorporación de funcionalidades más avanzadas sin comprometer la estabilidad de la operación actual.

\subsection{Alcance organizacional}

La implementación impacta transversalmente en la estructura vertical de la empresa:

\begin{itemize}
    \item[a.] \textbf{Gerencia General}: acceso a tableros de control (dashboards) con métricas en tiempo real sobre ventas y nacionalidades.
    \item[b.] \textbf{Ventas y Marketing}: digitalización del proceso de cotización y seguimiento sistemático de leads.
    \item[c.] \textbf{Operaciones y Reservas}: gestión eficiente de itinerarios y almacenamiento centralizado de vouchers y reservas.
    \item[d.] \textbf{Administración}: centralización de documentos y control de la eficiencia operativa.
\end{itemize}

El sistema será utilizado inicialmente por la gerencia y el personal operativo clave (como el apoyo familiar mencionado en las sesiones), definiendo roles de Administrador y Usuario de Ventas.

Adicionalmente, el documento de Alcance contempla componentes transversales como la migración inicial de datos (importación de la base de clientes y contactos, registro de datos maestros de servicios turísticos y configuración de usuarios y permisos) y la parametrización detallada de los flujos comerciales y operativos (etapas del embudo de ventas y del tablero de reservas). Estos elementos forman parte del trabajo de implementación descrito en las secciones posteriores de este informe.
\\
Con ello, el alcance no solo define qué módulos se configuran, sino también cómo se preparan las condiciones organizacionales y de datos para que la herramienta sea útil desde el primer día: quiénes la usarán, qué información mínima debe estar cargada y qué tipo de vistas y etapas se habilitan para reflejar la realidad operativa de la agencia de viajes.

\subsection{Alcance temporal}

El proyecto se ejecuta de manera progresiva siguiendo un cronograma de fases (sprints).

\subsection{Cronograma y estado de implementación}

El proyecto se ejecutará bajo una metodología ágil para asegurar que, al no ser un desarrollo desde cero sino una configuración y adaptación de flujos, la puesta en marcha sea rápida y eficiente. A continuación se presenta el cronograma de fases:

\begin{table}[h]
    \centering
    \begin{tabularx}{\textwidth}{|l|l|X|}
        \hline
        \textbf{Fase} & \textbf{Actividad clave} & \textbf{Enfoque en Odoo} \\
        \hline
        Fase 1 & Cimientos del sistema & Levantamiento de procesos y configuración técnica inicial de la instancia de Odoo en la nube. \\
        \hline
        Fase 2 & Migración de datos & Limpieza e importación de la base de datos histórica de clientes y catálogo de servicios turísticos. \\
        \hline
        Fase 3 & Configuración CRM y Proyectos & Parametrización de las etapas del embudo de ventas y los tableros Kanban para itinerarios. \\
        \hline
        Fase 4 & Automatización operativa & Configuración de plantillas de correo, mensajes de WhatsApp y registro de pagos manuales. \\
        \hline
        Fase 5 & Capacitación intensiva & Entrenamiento funcional del personal en el uso integral del sistema. \\
        \hline
        Fase 6 & Acompañamiento y cierre & Marcha blanca, ajustes finales en reportes y entrega del informe de cierre. \\
        \hline
    \end{tabularx}
    \caption{Cronograma y estado de implementación del proyecto}
\end{table}

